\pdfoutput=1
\documentclass[11pt]{amsart}
\usepackage{amsmath,amssymb,amsthm,booktabs,longtable}
\setcounter{MaxMatrixCols}{20}
\usepackage[margin=1.15in]{geometry}
\usepackage[colorlinks=true,linkcolor=blue,citecolor=blue,urlcolor=blue]{hyperref}
\hypersetup{pdfauthor={Bernd Johannes Wuebben},
pdftitle={Local stability chambers for toric tangent bundles}}
\newtheorem{theorem}{Theorem}[section]
\newtheorem{proposition}[theorem]{Proposition}
\newtheorem{lemma}[theorem]{Lemma}
\newtheorem{corollary}[theorem]{Corollary}
\theoremstyle{definition}
\newtheorem{definition}[theorem]{Definition}
\theoremstyle{remark}
\newtheorem{remark}[theorem]{Remark}
\newtheorem{example}[theorem]{Example}
\newcommand{\PP}{\mathbb P}
\newcommand{\RR}{\mathbb R}
\newcommand{\CC}{\mathbb C}
\newcommand{\ZZ}{\mathbb Z}
\newcommand{\Amp}{\operatorname{Amp}}
\newcommand{\rk}{\operatorname{rk}}
\newcommand{\Span}{\operatorname{span}}
\newcommand{\Gr}{\operatorname{Gr}}
\newcommand{\id}{\operatorname{id}}
\title{Local stability chambers for toric tangent bundles}
\author{Bernd Johannes Wuebben}
\subjclass[2020]{14M25 (primary); 14J45, 14J60 (secondary)}
\keywords{Toric variety, tangent bundle, slope stability, polarization,
projective bundle, Jordan--H\"older filtration}
\date{October 26, 2026}

\begin{document}
\begin{abstract}
We classify the anticanonical slope stability of the tangent bundles of
$X_{p,q}=\PP_{(\PP^1)^{p+q}}(\mathcal O\oplus\mathcal O(1^p,-1^q))$,
where $p,q$ are nonnegative integers and $p+q\ge1$. The tangent bundle is stable when
$p=q$, strictly semistable when $|p-q|=1$, and unstable otherwise.
For $p=k$ and $q=k-1$, the germ of the semistable polarization locus at
the anticanonical class is real analytically equivalent to
$[0,\infty)^k\times\RR^k$, with stable locus
$(0,\infty)^k\times\RR^k$. In particular, the number of independent
local chamber boundaries is unbounded.

For a smooth projective toric variety whose tangent bundle is semistable
at a polarization $H_0$, we describe the local chamber at $H_0$ by the order ideals of a finite poset
determined by the equal-slope ray subspaces. If each such subspace is
spanned by the rays it contains in every maximal cone, nearby stable
polarizations exist precisely when the variety is not a nontrivial
toric product. Applied to the existing anticanonical census in dimensions
three through six, this criterion gives nearby stability in all 31
indecomposable strictly semistable cases, whose local chambers have
one, two, or three independent boundaries. We determine locally free,
pairwise nonisomorphic Jordan--H\"older factors both for these cases and
for the family $X_{k,k-1}$.
\end{abstract}
\maketitle

\section{Introduction}

Slope stability depends on the polarization. A tangent bundle that is
strictly semistable at the anticanonical class can become stable after a
small change of that class, whereas a direct-sum decomposition prevents
stability for every polarization. We ask how this distinction is
reflected in the local geometry of the stable locus on a smooth toric
variety.

Let $X$ be smooth and projective of dimension $n$. For an ample real
divisor class $H$ and a torsion-free sheaf $E$, our convention is
\[
 \mu_H(E)=\frac{c_1(E)H^{n-1}}{\rk E}.
\]
The stable and semistable loci below are subsets of $\Amp(X)_\RR$.
We use quotient projectivization throughout: $\PP(E)$ parametrizes
one-dimensional quotients of $E$.

Our first result concerns a family of $\PP^1$-bundles. In
$\mathcal O(1^p,-1^q)$ the notation means $p$ entries equal to $1$
and $q$ entries equal to $-1$.

\begin{theorem}\label{thm:family}
For nonnegative integers $p,q$ with $p+q\ge1$, put
\[
 X_{p,q}=\PP_{(\PP^1)^{p+q}}
       \bigl(\mathcal O\oplus\mathcal O(1^p,-1^q)\bigr).
\]
These varieties are smooth Fano, and $T_{X_{p,q}}$ is, with respect to
$-K_{X_{p,q}}$,
\[
 \begin{cases}
 \text{stable},&p=q,\\
 \text{strictly semistable},&|p-q|=1,\\
 \text{unstable},&|p-q|\ge2.
 \end{cases}
\]
For $X_k=X_{k,k-1}$, of dimension and Picard number $2k$, the germ of
the semistable locus at $-K_{X_k}$ is real analytically equivalent to
\[
 [0,\infty)^k\times\RR^k,
\]
and the stable locus corresponds to $(0,\infty)^k\times\RR^k$.
There is an exact sequence of vector bundles
\begin{equation}\label{eq:intro-extension}
 0\longrightarrow T_{X_k/(\PP^1)^k}\longrightarrow T_{X_k}
 \longrightarrow\bigoplus_{i=1}^k\pi_i^*\mathcal O_{\PP^1}(2)
 \longrightarrow0,
\end{equation}
where $\pi_i$ is projection onto the $i$th positive-twist factor.
The relative tangent bundle is stable of rank $k$ at $-K_{X_k}$.
All displayed factors have the same slope and are pairwise nonisomorphic,
and every component of the extension class is nonzero.
\end{theorem}

The $k$ boundary hypersurfaces in this statement are the irredundant
boundaries of the semistable chamber. There are $2^k-1$ proper nonzero
equal-slope subspaces. For $k\ge2$, their equality hypersurfaces include hypersurfaces
defined by sums of several boundary coordinates, so the entire equality arrangement is
larger than the $k$ transverse chamber boundaries.

The proof reduces all slopes to integrals on $[-1,1]$. The anticanonical
polytope has cross-sectional volume
$(2+t)^p(2-t)^q$. The sign of a single odd moment separates the stable,
semistable, and unstable cases. At $p=q+1$, all equality subspaces contain
the vertical direction and every negative-twist direction. This identifies
both the chamber and the relative tangent factor in
\eqref{eq:intro-extension}.

The family also illustrates a more general description. The equal-slope
ray subspaces of a semistable toric tangent bundle form a finite
distributive lattice. Its associated poset determines the active slope
inequalities. The essential geometric issue is whether changes of
polarization vary the resulting slope gaps independently. In
Theorem~\ref{thm:general}, a condition on each maximal cone implies this
independence: a hypothetical principal relation among the block divisors
is successively annihilated by positive relations among rays. The local
chamber then has the form of an order-ideal cone, and its stable part is
nonempty exactly when the variety is not a toric product.

Finally, Section~\ref{sec:census} applies this criterion to the existing
anticanonical census in dimensions three through six. Among its 479
strictly semistable cases, 31 are indecomposable as toric products.
Their degrees and complete equal-slope subspaces are reconstructed here
from the ray configurations. All satisfy the geometric condition, and
their local chambers are orthants with one, two, or three boundary
coordinates. The other 448 varieties are products and have no stable
tangent polarization. The imported census supplies the exhaustive list;
the new finite calculation determines its local chamber geometry.

\subsection*{Relation to previous work}
The starting point is Klyachko's description of equivariant bundles
\cite{Klyachko} and the toric slope criterion in the form of
Hering--Nill--S\"u\ss\ \cite{HNS}. The reduction to finitely many ray
subspaces and the polyhedral description in curve classes are already
established there. The pullback to the ample cone is nonlinear, and
ample-cone stability regions need not be globally convex or polyhedral.
Our assertions concern their local analytic form.

Reynolds studies tangent stability on Fano Bott manifolds in Chapter~4
of \cite{Reynolds}. His Lemmas~4.2.3--4.2.4 compute the vertical facets
of $\PP^1$-bundles over products of projective lines, and
Proposition~4.2.6 gives a bound on their volumes. This is the ambient
construction used in Theorem~\ref{thm:family}. Our calculation also treats
all nonvertical ray subspaces and determines the complete equality
structure in every dimension. The Hirzebruch-surface case $k=1$ is
classical; see \cite[Section~6]{BDGP}. The Bott-tower constructions in
\cite[Section~6]{DDK} concern other equivariant rank-two bundles.

Napame--Tipler \cite{NT} give stability criteria for pullbacks along
toric fibrations and applications to stable perturbations.
Clarke--Tipler \cite{CT} describe local chambers and convergence of
Hermitian Yang--Mills connections for simple bundles whose graded object
is locally free with abelian automorphism group. The locally free, pairwise
nonisomorphic factors obtained here give the required graded objects.
The tangent bundles of $X_k$ and the 31 indecomposable census examples
are simple because they are stable at a nearby polarization. Thus these
examples satisfy the geometric hypotheses of \cite{CT}.
The distributive-lattice argument itself is the standard tight-set
construction for submodular polyhedra; see \cite{Fujishige}.

\section{Projective bundles and the number of chamber boundaries}
\label{sec:family}

\subsection{The fan and anticanonical degrees}
Write $d=p+q$, $n=d+1$, and let $h_i$ be the hyperplane classes on the
base. Set $a_i=1$ for $i\le p$, $a_i=-1$ otherwise,
$A=\sum_i a_i h_i$, $S=\sum_i h_i$, and
$\xi=c_1(\mathcal O_{X_{p,q}}(1))$. Then
\begin{equation}\label{eq:chow}
 h_i^2=0,\qquad \xi^2=A\xi,\qquad
 D:=-K_{X_{p,q}}=2\xi+2S-A.
\end{equation}
In a lattice basis $e_1,\ldots,e_d,v$, the fan rays are
\begin{equation}\label{eq:family-rays}
 v,\ -v,\qquad e_i,\ -e_i+a_iv\quad(1\le i\le d).
\end{equation}
Every maximal cone chooses one ray from each pair, including the vertical
pair. These are the projective-bundle cones of \cite[Section~7.3]{CLS}.
The divisor relations give $D_{-v}=\xi$, $D_v=\xi-A$, and both divisors
of the $i$th base pair have class $h_i$.

The anticanonical polytope in the dual coordinates $(x_1,\ldots,x_d,t)$
is
\begin{equation}\label{eq:family-polytope}
 -1\le t\le1,\qquad -1\le x_i\le1+a_it.
\end{equation}
At $t=\pm1$, all interval lengths $2+a_it$ are positive. Its normal
cones are therefore precisely the cones just described, and their
primitive generators are lattice bases. In particular $D$ is ample.

For the integral slope calculation through the next subsection, divide
divisor degrees by $(n-1)!$.
Thus facet volumes use Lebesgue measure in the affine facet lattice.
Put
\begin{equation}\label{eq:general-integrals}
 \begin{split}
 w(t)&=(2+t)^p(2-t)^q,\\
 M&=\int_{-1}^1w(t)\,dt,\\
 I_+&=\int_{-1}^1\frac{w(t)}{2+t}\,dt\quad(p>0),\\
 I_-&=\int_{-1}^1\frac{w(t)}{2-t}\,dt\quad(q>0).
 \end{split}
\end{equation}
The two vertical facets have degrees $3^p,3^q$. Each ray of a positive
pair has degree $I_+$, and each ray of a negative pair has degree $I_-$.
For a sloping facet $x_i=1+a_it$, the other coordinates and $t$ are
integral coordinates for its affine lattice; no additional normal-length
factor occurs. The full polytope has volume $M$, so
\begin{equation}\label{eq:total-slope}
 \frac{D^n}{(n-1)!}=nM=3^p+3^q+2pI_++2qI_-,
 \qquad \frac{\mu_D(T_X)}{(n-1)!}=M.
\end{equation}
Terms for an absent sign are omitted. The vertical formulas agree with
\cite[Lemmas~4.2.3--4.2.4]{Reynolds}.

\subsection{The slope test}
For a subspace $W\subset N_\CC$, the saturated induced equivariant
subsheaf $\mathcal F_W\subset T_X$ has rank $\dim W$ and
\begin{equation}\label{eq:toric-c1}
 c_1(\mathcal F_W)=\sum_{u_\rho\in W}D_\rho.
\end{equation}
It suffices to test slope inequalities on subspaces spanned by rays
\cite[Proposition~1.2 and Remark~2.6]{HNS}. For the configuration
\eqref{eq:family-rays}, these subspaces have a particularly simple form.
If $v\notin W$, at most one ray from each base pair belongs to $W$.
Their images in the base lattice are independent, so their number is
$\dim W$. If $v\in W$, either both rays of a base pair belong to $W$
or neither does; hence
\begin{equation}\label{eq:vertical-flat}
 W=\Span_\CC(v,e_i:i\in T)
\end{equation}
for a subset $T\subseteq\{1,\ldots,d\}$.

Interchanging $p,q$ reverses the vertical coordinate, so suppose $p\ge q$.
If $p=q\ge1$, symmetry gives $I_+=I_-=I$. We have $M>I$ because
$2+t>1$ in the interior, whereas
\[
 M-2I=\int_{-1}^1t(2-t)(4-t^2)^{p-1}\,dt
      =-\int_{-1}^1t^2(4-t^2)^{p-1}\,dt<0.
\]
A subspace not containing $v$ thus has normalized slope $I<M$.
For a proper subspace containing $v$, let $J$ be the nonempty set of
omitted base indices. By \eqref{eq:total-slope}, its normalized slope
margin is
\[
 \frac{c_1(\mathcal F_W)D^{n-1}}{(n-1)!}-(\dim W)M
       =|J|(M-2I)<0.
\]
This proves stability in the balanced case.

If $p\ge q+2$, write $h=p-q-1\ge1$. Then
\[
 M-2I_+=\int_{-1}^1t(2+t)^h(4-t^2)^q\,dt>0,
\]
since pairing $t$ and $-t$ gives the positive integrand
$t((2+t)^h-(2-t)^h)(4-t^2)^q$ for $0<t<1$.
Take $W$ to contain $v$ and every base direction except one positive
direction. Its rank is $n-1$ and its normalized slope margin is
$M-2I_+>0$. Thus the tangent bundle is unstable.

It remains to treat $p=q+1$. Put $k=p$ and $q=k-1$. Define
\[
 I_q=\int_{-1}^1(4-t^2)^q\,dt,
\]
and, for $q\ge1$, put
\[
 J_q=\int_{-1}^1(4+t^2)(4-t^2)^{q-1}\,dt.
\]
These are $I_+$ and $I_-$, respectively, and $M=2I_q$. Integrating
the derivative of $t(4-t^2)^q$ gives
\begin{equation}\label{eq:integral-identity}
 qJ_q=(q+1)I_q-2\cdot3^q\qquad(q\ge1).
\end{equation}
For $q\ge1$ we have $I_q>2\cdot3^q$, and it follows that
\begin{equation}\label{eq:integral-bounds}
 I_q<J_q<2I_q\qquad(q\ge1).
\end{equation}
Thus every subspace not containing $v$ has slope strictly less than
$M$. For $q\ge1$, a subspace of the form \eqref{eq:vertical-flat}, with $a$ positive
and $b$ negative base directions, has rank $1+a+b$ and margin
\begin{equation}\label{eq:family-margin}
 4\cdot3^q+2aI_q+2bJ_q-2(1+a+b)I_q
       =2(b-q)(J_q-I_q).
\end{equation}
For $q\ge1$ this is nonpositive, with equality exactly when $b=q$.
For $q=0$, there are no negative directions and $I_0=2$; the margin is
directly $4+2aI_0-2(1+a)I_0=0$.

Consequently the proper nonzero equal-slope subspaces are precisely
\begin{equation}\label{eq:family-equalities}
 V_0+\Span_\CC(e_i:i\in T),\qquad
 V_0=\Span_\CC(v,e_{k+1},\ldots,e_{2k-1}),\quad
 T\subsetneq\{1,\ldots,k\}.
\end{equation}
This proves the trichotomy and exhausts all equalities.

\subsection{Local coordinates and stable factors}
We now return to ordinary intersection degrees, without the factorial
normalization used in the integral calculation.
The space $V_0$ in \eqref{eq:family-equalities} corresponds to the relative
tangent bundle $F=T_{X_k/(\PP^1)^k}$. Indeed the projection to the
positive base factors is locally a product with fiber
\[
 \PP_{(\PP^1)^{k-1}}
       \bigl(\mathcal O\oplus\mathcal O(-1,\ldots,-1)\bigr).
\]
Its relative tangent sequence is \eqref{eq:intro-extension}.
Write $C_0=c_1(F)=D-2\sum_{i=1}^kh_i$ and $C_i=2h_i$ for $i\ge1$.
These classes are independent: $C_0$ has nonzero $\xi$ coefficient in
the basis $\xi,h_1,\ldots,h_{2k-1}$.

For a nearby polarization set
\begin{equation}\label{eq:family-z}
 z_i(H)=\left(2h_i-\frac{D}{2k}\right)H^{2k-1},\qquad1\le i\le k.
\end{equation}
The centered divisor classes in this expression have rank $k$.
The derivative of $H\mapsto H^{n-1}$ at an ample class is an isomorphism
on numerical divisor and curve spaces, by the Hodge-index pairing
\cite[Theorem~6.4 and Corollary~6.5]{GRT}. Thus the $z_i$ extend to
real analytic local coordinates at $D$.

Only the inequalities in \eqref{eq:family-equalities} can be active
sufficiently near $D$, since the finite list of other inequalities is
strict there. For the subspace indexed by $T$, its slope margin is
$-\sum_{i\notin T}z_i(H)$. Hence semistability is exactly $z_i\ge0$
for every $i$, and stability exactly $z_i>0$ for every $i$. This proves
the chamber statement of Theorem~\ref{thm:family}.

The subbundle $F$ is stable at $D$. A subsheaf of greater slope would
contradict semistability of $T_X$; an equivariant equal-slope proper
subsheaf would contradict the minimality of $V_0$ in
\eqref{eq:family-equalities}. The toric criterion tests ordinary
stability, so this suffices. The quotient factors are line bundles of
the same slope. Their first Chern classes, together with $C_0$, are
independent, which excludes isomorphic factors. Finally a zero
component of the extension class would split off its quotient line
bundle. A nontrivial direct sum cannot be stable for any polarization,
whereas the local chamber contains stable classes. Every extension
component is therefore nonzero, completing the proof of
Theorem~\ref{thm:family}.

\begin{example}
For $k=1$, $X_k$ is the Hirzebruch surface $\mathbb F_1$. If $S$ is
its negative section and $f$ a fiber, the first Chern classes of the two
stable factors are
$2S+f$ and $2f$. For $H=aS+bf$ with $b>a>0$, the relevant inequality
is $2b-3a<0$. Thus the stable side near $-K=2S+3f$ is
$a<b<3a/2$, in agreement with the classical surface calculation.
For $k=2,3$, the anticanonical degrees are $352$ and $38976$,
respectively; the proper equality subspaces number $3$ and $7$.
\end{example}

\section{Equal-slope subspaces and local chambers}
\label{sec:general}

Let $X=X_\Sigma$ be a smooth projective toric variety of dimension
$n\ge2$. Write $R=\Sigma(1)$, let $u_\rho$ be the primitive ray
generators, and put $D=\sum_{\rho\in R}D_\rho=c_1(T_X)$.
Fix an ample real class $H_0$ such that $T_X$ is semistable, and set
\[
 \alpha_0=H_0^{n-1},\qquad d_\rho=D_\rho\alpha_0>0,
 \qquad\mu=D\alpha_0/n>0.
\]
Positivity holds because each invariant divisor is nonzero and effective.
For a subspace $V\subset N_\CC$ write
$R_V=\{\rho:u_\rho\in V\}$ and $d(V)=\sum_{\rho\in R_V}d_\rho$.

We use the toric criterion \eqref{eq:toric-c1} also for real ample
classes. To justify this extension, specialize an arbitrary subsheaf
to a torus-fixed point in the closure of its orbit in the projective
Quot scheme. Flat specialization preserves rank and numerical first
Chern class, and saturation only increases degree against an ample
real class. Thus any ordinary destabilizing or equal-slope subsheaf
gives an equivariant one with at least that slope. This is the same
equivariant reduction underlying \cite[Proposition~2.3]{HNS}.

Let $\mathcal L$ consist of $0,N_\CC$ and the subspaces with
$d(V)=\mu\dim V$. Each such subspace is ray-spanned: otherwise its
contained-ray span has the same positive degree and smaller rank,
contradicting semistability. Thus $\mathcal L$ is finite.

\begin{lemma}\label{lem:lattice}
The set $\mathcal L$ is a distributive lattice under sum and
intersection. For $V,W\in\mathcal L$,
\[
 R_{V+W}=R_V\cup R_W,\qquad R_{V\cap W}=R_V\cap R_W.
\]
\end{lemma}
\begin{proof}
Ray membership and semistability give
\[
 d(V)+d(W)\le d(V+W)+d(V\cap W)
 \le\mu\bigl(\dim(V+W)+\dim(V\cap W)\bigr)
 =d(V)+d(W).
\]
Both inequalities are equalities. The two resulting subspaces have
the required slope, and positivity of all $d_\rho$ excludes any extra
ray in $V+W$. Since equality subspaces are ray-spanned, $V\mapsto R_V$
is injective. Its image is closed under union and intersection, which
proves distributivity.
\end{proof}

This is the standard tight-set argument for a submodular rank function
\cite{Fujishige}. In fact $(d_\rho/\mu)_{\rho\in R}$ belongs to the
base polytope of the ray matroid, and its tight sets are exactly the
sets $R_V$ above.

Partition $R$ into blocks $B_p$ according to identical membership in
all equality sets. Define $p\le q$ if every equality set containing
$B_q$ also contains $B_p$, and denote this finite poset by $P$.
The equality sets are precisely unions of blocks over order ideals of
$P$. Indeed, intersection of all equality sets containing $B_q$ gives
its principal ideal; unions of these give every ideal.
For an ideal $I$, denote its subspace by $V_I$.

Put $s=|P|$ and $C_p=\sum_{\rho\in B_p}D_\rho$. There are positive
integers $r_p$ such that
\begin{equation}\label{eq:rank-valuation}
 \dim V_I=\sum_{p\in I}r_p,\qquad
 \sum_pr_p=n,\qquad C_p\alpha_0=\mu r_p.
\end{equation}
To see this, choose a linear extension of $P$. The rank increment
when $B_p$ is added to an initial ideal equals
$\sum_{\rho\in B_p}d_\rho/\mu$. It is positive and integral, and the
formula shows independence of the chosen linear extension.

\begin{definition}
In $Z=\{z\in\RR^s:\sum_pz_p=0\}$ define
\[
 \mathcal C(P)=\left\{z:\sum_{p\in I}z_p\le0
       \text{ for every nonempty proper order ideal }I\right\}.
\]
Let $\mathcal C^\circ(P)$ denote the set where all displayed
inequalities are strict. The superscript denotes this strict locus;
it does not assert that the cone is full-dimensional.
\end{definition}

We impose the following condition on equal-slope subspaces:
\begin{equation}\label{eq:cone-condition}
 V=\Span_\CC\{u_\rho:\rho\in\sigma(1),\ u_\rho\in V\}
 \quad\text{for every }V\in\mathcal L
 \text{ and every maximal cone }\sigma.
\end{equation}
Because the rays of a maximal smooth cone form a basis, this says that
exactly $\dim V$ of them lie in $V$.

\begin{theorem}\label{thm:general}
Under \eqref{eq:cone-condition}, the germs of the semistable and stable
polarization loci at $H_0$ are real analytically equivalent to
\[
 \mathcal C(P)\times\RR^{\rho(X)-s+1},\qquad
 \mathcal C^\circ(P)\times\RR^{\rho(X)-s+1},
\]
respectively. Stable polarizations occur arbitrarily close to $H_0$
if and only if $X$ is not a nontrivial toric product. In that case
rational ample stable classes occur arbitrarily close to $H_0$.
\end{theorem}

\begin{lemma}\label{lem:positive-rays}
If a ray subspace $V$ satisfies \eqref{eq:cone-condition}, its induced
sheaf $\mathcal F_V$ is a vector subbundle of $T_X$. The cones of
$\Sigma$ contained in $V_\RR$ form a complete fan in $V_\RR$.
\end{lemma}
\begin{proof}
On each smooth affine chart, $T_X$ splits into the coordinate weight
lines of its ray basis. The condition makes $V$ a sum of a subset of
these lines, so the induced inclusion is locally a direct summand.
For completeness, take $v\in V_\RR$ and a maximal cone containing it.
In that ray basis, all coordinates outside $V$ vanish. Thus $v$ lies
in the face generated by the rays in $V$. These faces cover $V_\RR$.
\end{proof}

\begin{lemma}\label{lem:independence}
Under \eqref{eq:cone-condition}, the numerical classes $C_p$ are
linearly independent.
\end{lemma}
\begin{proof}
Numerical and linear equivalence agree for divisors on a smooth
complete toric variety, and the invariant divisor sequence identifies
$N^1(X)_\RR$ with $\RR^R/M_\RR$
\cite[Propositions~6.3.15 and 6.4.1]{CLS}.
Consequently a relation $\sum_pa_pC_p=0$ is represented by a character
$m\in M_\RR$ satisfying $\langle m,u_\rho\rangle=a_p$ on $B_p$.

Order the blocks by a linear extension $p_1,\ldots,p_s$, and let $I_j$
be the initial ideal. By Lemma~\ref{lem:positive-rays}, its rays
positively span $V_{I_j}$. They therefore admit a relation
$\sum_{\rho\in R_{V_{I_j}}}b_\rho u_\rho=0$ with every $b_\rho>0$:
sum nonnegative expressions for the negatives of all generators.
Inductively suppose $a_{p_i}=0$ for $i<j$. Pairing this positive
relation with $m$ yields
\[
 a_{p_j}\sum_{\rho\in B_{p_j}}b_\rho=0,
\]
and hence $a_{p_j}=0$. All coefficients vanish.
\end{proof}

\begin{proof}[Proof of Theorem~\ref{thm:general}]
For a curve class $\alpha$ near $\alpha_0$, set
\[
 z_p(\alpha)=\left(C_p-\frac{r_p}{n}D\right)\alpha.
\]
These coordinates sum to zero and vanish at $\alpha_0$.
Lemma~\ref{lem:independence} implies that their linear map onto $Z$
has rank $s-1$. The Hodge-index argument used in
Section~\ref{sec:family} shows that $H\mapsto z(H^{n-1})$ is a
submersion at $H_0$. Complete its components to real analytic local
coordinates. All initially strict inequalities for ray-spanned subspaces remain strict
in a sufficiently small neighborhood. The remaining inequalities are
exactly those defining $\mathcal C(P)$, by
\eqref{eq:rank-valuation}.

The strict cone is nonempty exactly when the undirected Hasse diagram
of $P$ is connected. If it is disconnected, a connected component and
its complement are both ideals, whose coordinate sums cannot both be
negative. If it is connected, orient each covering relation upward
and sum $e_q-e_p$ over the edges $p<q$. For every nonempty proper
ideal, no edge enters and at least one leaves, so its coordinate sum
is negative. For a singleton poset the condition is vacuous.

It remains to interpret connectedness. Disconnectedness gives two
complementary equality ideals. Their ray spans have ranks summing to
$n$, and every ray belongs to one of them, so they split $N_\RR$.
A maximal-cone lattice basis shows that this is an integral lattice
splitting. Every cone splits into its component faces. Completeness
forces every combination of component cones to occur: take interior
points in the two cones and a cone containing their sum; the fan
intersection property makes their sum a face of that cone.
Thus the fan is a product fan. Conversely, a toric product splits
$T_X$ into equivariant summands. Semistability and the weighted-average
identity force all summands to have slope $\mu$. Their complementary
ray sets are equality ideals, which disconnect $P$.

Finally, strict inequalities define an open subset. Rational ample
classes are dense there. A nontrivial direct sum cannot be stable at
any polarization, since one summand has slope at least the total
weighted average.
\end{proof}

\begin{remark}
Only rank $s-1$ of the centered divisor classes is needed for the
coordinate argument. Condition \eqref{eq:cone-condition} is a sufficient
geometric condition that implies the stronger independence of all
$C_p$. We do not assert that every semistable toric tangent bundle
satisfies it.
\end{remark}

\subsection{The graded bundle}
\begin{proposition}\label{prop:graded}
Under \eqref{eq:cone-condition}, a maximal chain of equality subspaces
gives a Jordan--H\"older filtration by vector subbundles. Its stable
quotients have ranks $r_p$ and first Chern classes $C_p$, and are
pairwise nonisomorphic. In particular the graded bundle is locally
free with automorphism group $(\CC^*)^s$.
\end{proposition}
\begin{proof}
Local freeness follows from Lemma~\ref{lem:positive-rays}; additivity
gives the ranks and Chern classes. Write $A\subset B$ for successive
terms. A saturated equivariant subsheaf $G\subset B/A$ lifts to a
subsheaf $\widetilde G\subset T_X$. It is saturated, since
\[
 0\longrightarrow(B/A)/G\longrightarrow T_X/\widetilde G
 \longrightarrow T_X/B\longrightarrow0
\]
has torsion-free end terms. If $G$ had greater slope, so would
$\widetilde G$, contrary to semistability. If it had equal slope,
its subspace would lie strictly between two successive terms of the
maximal equality chain. This is impossible. The toric criterion
therefore gives ordinary stability. Independence of $C_p$ excludes
isomorphic quotients. Stable bundles of equal slope have only scalar
endomorphisms and no nonzero morphisms between distinct factors.
\end{proof}

\section{The strictly semistable census in dimensions three through six}
\label{sec:census}

The anticanonical census of \cite{Wuebben} contains 479 strictly
semistable cases in dimensions three through six. Its full scope is
imported here; the present calculation reconstructs the local data of
the indecomposable cases. The identifiers below refer to its smooth
Fano polytope lists, with primitive fan generators as vertices.

\begin{theorem}\label{thm:census}
Among the strictly semistable cases of the dimension three through six
census, the 31 varieties that are not nontrivial toric products admit
stable polarizations arbitrarily close to the anticanonical class.
Each satisfies \eqref{eq:cone-condition}. Its equality poset has one
minimal element and $k$ incomparable elements above it, with counts
\[
\begin{array}{c|rrr|rr}
 n&k=1&k=2&k=3&\text{indecomposable}&\text{product}\\ \hline
 3&2&0&0&2&5\\
 4&2&1&0&3&21\\
 5&7&1&0&8&74\\
 6&15&2&1&18&348
\end{array}
\]
The local semistable germ is
$[0,\infty)^k\times\RR^{\rho(X)-k}$, with stable locus its strict
orthant. The other 448 cases have no stable tangent polarization.
\end{theorem}

\begin{proof}
For each selected indecomposable configuration, the calculation gives
the following complete description. There is a minimal equal-slope
ray block $B_0$, of rank $r_0$, and upper blocks $B_1,\ldots,B_k$.
The nonzero proper equality sets are exactly $B_0$ together with any
proper subset of the upper blocks. Their ranks are additive, and
every equality subspace meets every maximal cone in exactly its rank
many ray generators. The ray matrices, degree vectors, and block
partitions are printed in Appendix~\ref{app:data}; the summary appears
in Table~\ref{tab:records}. Theorem~\ref{thm:general} applies.

For this poset, the ideal containing all blocks except $B_j$ gives
$z_j\ge0$. These $k$ inequalities imply all the other ideal
inequalities, and their strict versions do the same for stability.
Thus the local cone is an orthant. The remaining selected ray
configurations split into nontrivial direct sums. Smoothness and
completeness convert these into toric products, as in the proof of
Theorem~\ref{thm:general}. Their tangent bundles cannot be stable.
\end{proof}

The number $k$ need not equal the rank of the quotient by the minimal
subbundle. For example, F.6D.6547 has block ranks $(3,3)$ and a
single chamber boundary. The multiple-boundary cases have ranks
$(2,1,1)$, $(3,1,1)$, $(4,1,1)$, $(4,1,1)$, and $(3,1,1,1)$.
The first and last are $X_2$ and $X_3$ from Theorem~\ref{thm:family}.
In the coordinates of Appendix~\ref{app:data}, the images of the ordered
family basis $(e_1,\ldots,e_{2k-1},v)$ are
\[
 (-e_1,-e_3,-e_2,e_4)\quad(k=2),\qquad
 (-e_1,-e_2,-e_5,-e_3,-e_4,e_6)\quad(k=3).
\]
These signed permutations identify the full fans, not only their ray sets.

In all 31 cases, let $F$ be the stable minimal subbundle. The images
of the upper blocks in $N_\CC/V_0$ span a direct sum: they span the
quotient and their ranks add to its dimension. Each projected ray
filtration belongs to one summand, so the quotient bundle splits.
We obtain
\begin{equation}\label{eq:census-extension}
 0\longrightarrow F\longrightarrow T_X
 \longrightarrow Q_1\oplus\cdots\oplus Q_k\longrightarrow0,
\end{equation}
with stable locally free factors of the same slope, pairwise
nonisomorphic by Proposition~\ref{prop:graded}. Each extension
component is nonzero, by nearby stability and the splitting argument
used for \eqref{eq:intro-extension}. Thus 26 cases give nonsplit
extensions of two stable bundles; the remaining five give three or
four stable factors. The product cases need not be polystable:
a product can have a factor with strictly semistable nonpolystable
tangent bundle.

% Generated from the exact chamber data.
\begin{longtable}{lrrrrl}
\caption{The indecomposable strictly semistable cases. The last column gives the block ranks.}\label{tab:records}\\
\toprule Identifier & $n$ & $\rho$ & $k$ & $\mu_{-K}(T_X)$ & Ranks\\\midrule\endfirsthead
\toprule Identifier & $n$ & $\rho$ & $k$ & $\mu_{-K}(T_X)$ & Ranks\\\midrule\endhead
\bottomrule\endfoot
F.3D.0003 & 3 & 5 & 1 & $12$ & $(2,1)$\\
F.3D.0013 & 3 & 2 & 1 & $18$ & $(2,1)$\\
F.4D.0072 & 4 & 4 & 2 & $88$ & $(2,1,1)$\\
F.4D.0084 & 4 & 4 & 1 & $88$ & $(3,1)$\\
F.4D.0114 & 4 & 2 & 1 & $128$ & $(3,1)$\\
F.5D.0245 & 5 & 7 & 1 & $460$ & $(4,1)$\\
F.5D.0320 & 5 & 6 & 1 & $498$ & $(4,1)$\\
F.5D.0391 & 5 & 6 & 1 & $498$ & $(4,1)$\\
F.5D.0409 & 5 & 5 & 1 & $652$ & $(4,1)$\\
F.5D.0496 & 5 & 7 & 2 & $536$ & $(3,1,1)$\\
F.5D.0536 & 5 & 7 & 1 & $536$ & $(4,1)$\\
F.5D.0569 & 5 & 7 & 1 & $536$ & $(4,1)$\\
F.5D.0841 & 5 & 2 & 1 & $1250$ & $(4,1)$\\
F.6D.1739 & 6 & 4 & 1 & $7224$ & $(5,1)$\\
F.6D.2022 & 6 & 4 & 1 & $7224$ & $(5,1)$\\
F.6D.3068 & 6 & 6 & 1 & $5832$ & $(5,1)$\\
F.6D.3115 & 6 & 6 & 1 & $5848$ & $(5,1)$\\
F.6D.3127 & 6 & 6 & 1 & $5848$ & $(5,1)$\\
F.6D.3996 & 6 & 5 & 1 & $7404$ & $(5,1)$\\
F.6D.4188 & 6 & 6 & 1 & $6376$ & $(5,1)$\\
F.6D.4954 & 6 & 5 & 1 & $7452$ & $(5,1)$\\
F.6D.5485 & 6 & 8 & 2 & $5160$ & $(4,1,1)$\\
F.6D.5541 & 6 & 6 & 3 & $6496$ & $(3,1,1,1)$\\
F.6D.5880 & 6 & 8 & 1 & $5160$ & $(5,1)$\\
F.6D.5962 & 6 & 6 & 1 & $6496$ & $(5,1)$\\
F.6D.5965 & 6 & 5 & 2 & $7236$ & $(4,1,1)$\\
F.6D.6078 & 6 & 8 & 1 & $5160$ & $(5,1)$\\
F.6D.6547 & 6 & 4 & 1 & $\frac{25984}{3}$ & $(3,3)$\\
F.6D.6614 & 6 & 5 & 1 & $7236$ & $(5,1)$\\
F.6D.7472 & 6 & 4 & 1 & $9024$ & $(5,1)$\\
F.6D.7551 & 6 & 2 & 1 & $15552$ & $(5,1)$\\
\end{longtable}


\subsection{Exact computation and its scope}
The fan of a smooth Fano polytope consists of cones over its faces.
For the 31 configurations we determine every facet by its rational
supporting equation and check that its vertices form an integral basis.
Every other ray lies strictly below its supporting hyperplane, and
every ridge occurs in two facets. This nonempty collection is closed
under facet adjacency and hence is the entire boundary. An independent
calculation enumerates all subsets of $n$ rays and recovers the same
facets directly.

For degrees we use equivariant localization. If $m_{\sigma,i}$ is
the character basis dual to the rays of a maximal cone and
$w_{\sigma,i}=\langle m_{\sigma,i},\eta\rangle$ for a generic
integral $\eta$, then
\[
 D^n=\sum_\sigma\frac{(\sum_iw_{\sigma,i})^n}{\prod_iw_{\sigma,i}},
 \quad
 D_\rho D^{n-1}=
 \sum_{\sigma\ni\rho}
 \frac{w_{\sigma,\rho}(\sum_iw_{\sigma,i})^{n-1}}
      {\prod_iw_{\sigma,i}}.
\]
The same integers are independently obtained by triangulating the
corresponding facets of the polar polytope and summing their lattice
volumes. Ray flats are enumerated by adjoining a ray and taking linear
closure; every flat has a ray basis, which proves completeness of this
procedure. Independent enumeration by intersections of ray-spanned
hyperplanes gives the same 4,890 proper flats. Every slope inequality
and every equality is then tested exactly.

The product partitions of all 479 selected cases are obtained from
the connected components of their linear ray matroids. Fundamental
circuits relative to one basis and exhaustive minimal dependent
subsets give identical partitions. The complete census and its
exclusion of other strictly semistable cases remain inputs from
\cite{Wuebben}; they are not consequences of checking only the 31
indecomposable records. The accompanying data and rational-arithmetic
programs specify the input hashes and reproduce the calculations.
The supplementary directory \texttt{reproducibility/} contains the
31 complete records, the rays of all 479 selected cases, and an exact
verifier requiring only Python's standard library. Its \texttt{README.md}
gives the command \texttt{python3 verify.py}, the data conventions, and
the optional comparison with the full source files. This verifier
reconstructs facets, degrees, flats, and product partitions; selection
of the 479 cases from the complete census remains an imported input.

\appendix
\section{Ray configurations, degrees, and equality blocks}\label{app:data}
For each record, the columns of $U$ are the primitive ray generators
in their given order, numbered from $0$. The vector $d$ lists
$D_\rho(-K_X)^{n-1}$. The blocks $B_0;B_1;\ldots;B_k$ determine
every proper equality subspace as described in the proof of
Theorem~\ref{thm:census}. The fan is recovered by taking cones over
the faces of $\operatorname{conv}(U)$, so the displayed data determine
all inputs to the 31 indecomposable computations.
% Generated from the exact chamber data.
\par\medskip\noindent\begin{minipage}{\linewidth}\textbf{F.3D.0003.}
\begingroup\small\setlength{\arraycolsep}{3pt}
\[U=\begin{pmatrix}
-1 & -1 & -1 & 0 & 0 & 0 & 1 & 1\\
1 & 0 & 0 & -1 & 0 & 0 & 0 & 0\\
1 & 0 & 1 & 0 & 1 & -1 & 0 & -1\\
\end{pmatrix},\qquad \mu=12.\]
\[d=(6,3,2,6,3,5,5,6).\]
\[B_0=\{1,2,4,5,6,7\};\quad B_1=\{0,3\}.\]
\endgroup\end{minipage}\par
\par\medskip\noindent\begin{minipage}{\linewidth}\textbf{F.3D.0013.}
\begingroup\small\setlength{\arraycolsep}{3pt}
\[U=\begin{pmatrix}
-1 & 1 & -1 & 0 & 0\\
1 & 0 & 0 & -1 & 0\\
0 & 1 & 0 & 0 & -1\\
\end{pmatrix},\qquad \mu=18.\]
\[d=(9,15,6,9,15).\]
\[B_0=\{1,2,4\};\quad B_1=\{0,3\}.\]
\endgroup\end{minipage}\par
\par\medskip\noindent\begin{minipage}{\linewidth}\textbf{F.4D.0072.}
\begingroup\small\setlength{\arraycolsep}{3pt}
\[U=\begin{pmatrix}
0 & 0 & 0 & 0 & -1 & 1 & 0 & 0\\
0 & 0 & -1 & 0 & 0 & 0 & 0 & 1\\
1 & -1 & 0 & 0 & 0 & 0 & 0 & 0\\
1 & 0 & 0 & -1 & 0 & 1 & 1 & -1\\
\end{pmatrix},\qquad \mu=88.\]
\[d=(44,44,52,54,44,44,18,52).\]
\[B_0=\{2,3,6,7\};\quad B_1=\{0,1\};\quad B_2=\{4,5\}.\]
\endgroup\end{minipage}\par
\par\medskip\noindent\begin{minipage}{\linewidth}\textbf{F.4D.0084.}
\begingroup\small\setlength{\arraycolsep}{3pt}
\[U=\begin{pmatrix}
0 & 0 & 0 & 0 & -1 & 0 & 0 & 1\\
1 & 0 & 0 & -1 & 0 & 1 & 1 & -1\\
0 & 0 & -1 & 0 & 0 & 0 & 1 & 0\\
0 & -1 & 0 & 0 & 0 & 1 & 1 & 0\\
\end{pmatrix},\qquad \mu=88.\]
\[d=(27,68,44,45,50,24,44,50).\]
\[B_0=\{0,1,3,4,5,7\};\quad B_1=\{2,6\}.\]
\endgroup\end{minipage}\par
\par\medskip\noindent\begin{minipage}{\linewidth}\textbf{F.4D.0114.}
\begingroup\small\setlength{\arraycolsep}{3pt}
\[U=\begin{pmatrix}
0 & -1 & 1 & 0 & 0 & 0\\
0 & 0 & 1 & -1 & 0 & -1\\
0 & 0 & 0 & 0 & -1 & 1\\
-1 & 0 & 1 & 0 & 0 & 0\\
\end{pmatrix},\qquad \mu=128.\]
\[d=(112,112,112,48,64,64).\]
\[B_0=\{0,1,2,3\};\quad B_1=\{4,5\}.\]
\endgroup\end{minipage}\par
\par\medskip\noindent\begin{minipage}{\linewidth}\textbf{F.5D.0245.}
\begingroup\small\setlength{\arraycolsep}{3pt}
\[U=\begin{pmatrix}
-1 & -1 & 0 & 0 & 1 & 0 & 0 & 0 & 0 & 0 & 0 & 1\\
0 & -1 & 0 & 0 & 0 & 0 & -1 & 0 & 0 & 1 & 0 & 1\\
0 & 0 & 1 & 0 & 0 & 0 & 0 & 0 & 0 & 0 & -1 & 0\\
0 & 1 & 0 & 1 & 0 & 0 & 0 & -1 & 0 & 0 & 0 & -1\\
0 & 1 & -1 & 0 & 0 & -1 & 0 & 0 & 1 & 0 & 0 & -1\\
\end{pmatrix},\qquad \mu=460.\]
\[d=(161,207,230,161,207,92,161,207,276,207,230,161).\]
\[B_0=\{0,1,3,4,5,6,7,8,9,11\};\quad B_1=\{2,10\}.\]
\endgroup\end{minipage}\par
\par\medskip\noindent\begin{minipage}{\linewidth}\textbf{F.5D.0320.}
\begingroup\small\setlength{\arraycolsep}{3pt}
\[U=\begin{pmatrix}
0 & -1 & -1 & 0 & 1 & 0 & 1 & 0 & 0 & 0 & 1\\
0 & 0 & 0 & 0 & 1 & -1 & 0 & 0 & 0 & 0 & 0\\
1 & 0 & 1 & 0 & -1 & 0 & -1 & 0 & -1 & 1 & -1\\
1 & 0 & 0 & 0 & 0 & 0 & 0 & -1 & 0 & 1 & -1\\
0 & 0 & 0 & -1 & 1 & 0 & 1 & 0 & 0 & -1 & 1\\
\end{pmatrix},\qquad \mu=498.\]
\[d=(177,279,260,260,249,249,104,270,177,279,186).\]
\[B_0=\{0,1,2,3,6,7,8,9,10\};\quad B_1=\{4,5\}.\]
\endgroup\end{minipage}\par
\par\medskip\noindent\begin{minipage}{\linewidth}\textbf{F.5D.0391.}
\begingroup\small\setlength{\arraycolsep}{3pt}
\[U=\begin{pmatrix}
0 & -1 & 0 & 0 & 0 & 0 & 1 & -1 & 1 & 0 & -1\\
0 & 0 & 0 & -1 & 0 & 0 & 0 & 1 & 0 & 0 & 0\\
0 & 1 & 1 & 0 & -1 & 0 & -1 & 0 & -1 & 1 & 0\\
-1 & 0 & 1 & 0 & 0 & 0 & 0 & 0 & -1 & 1 & 0\\
0 & 0 & -1 & 0 & 0 & -1 & 1 & 0 & 1 & 0 & 0\\
\end{pmatrix},\qquad \mu=498.\]
\[d=(229,166,310,249,166,229,229,249,310,229,124).\]
\[B_0=\{0,1,2,4,5,6,8,9,10\};\quad B_1=\{3,7\}.\]
\endgroup\end{minipage}\par
\par\medskip\noindent\begin{minipage}{\linewidth}\textbf{F.5D.0409.}
\begingroup\small\setlength{\arraycolsep}{3pt}
\[U=\begin{pmatrix}
0 & 0 & 0 & -1 & 1 & 0 & 0 & -1 & -1 & 1\\
-1 & 0 & 0 & 1 & 0 & 0 & 0 & 0 & 0 & 0\\
0 & 1 & 0 & 0 & 0 & -1 & 0 & 0 & 0 & 0\\
0 & 0 & -1 & 0 & 0 & 0 & 0 & 1 & 0 & 0\\
0 & 1 & 0 & 1 & 0 & 0 & -1 & 0 & 0 & -1\\
\end{pmatrix},\qquad \mu=652.\]
\[d=(342,378,326,342,336,378,276,326,112,444).\]
\[B_0=\{0,1,3,4,5,6,8,9\};\quad B_1=\{2,7\}.\]
\endgroup\end{minipage}\par
\par\medskip\noindent\begin{minipage}{\linewidth}\textbf{F.5D.0496.}
\begingroup\small\setlength{\arraycolsep}{3pt}
\[U=\begin{pmatrix}
0 & 0 & 0 & 0 & 0 & 0 & -1 & 0 & 0 & 0 & 0 & 1\\
0 & -1 & 0 & 0 & 1 & 0 & 0 & -1 & -1 & 1 & 1 & -1\\
0 & 0 & 0 & -1 & 0 & 0 & 0 & 0 & 0 & 0 & 1 & 0\\
-1 & 1 & 0 & 0 & 0 & 0 & 0 & 0 & 0 & 0 & 0 & 0\\
0 & 1 & 1 & 0 & 0 & -1 & 0 & 0 & 1 & -1 & -1 & 1\\
\end{pmatrix},\qquad \mu=536.\]
\[d=(268,268,134,308,218,218,268,134,72,216,308,268).\]
\[B_0=\{2,3,4,5,7,8,9,10\};\quad B_1=\{0,1\};\quad B_2=\{6,11\}.\]
\endgroup\end{minipage}\par
\par\medskip\noindent\begin{minipage}{\linewidth}\textbf{F.5D.0536.}
\begingroup\small\setlength{\arraycolsep}{3pt}
\[U=\begin{pmatrix}
-1 & -1 & -1 & 0 & 0 & 0 & 0 & 0 & 0 & 1 & 1 & 1\\
0 & 0 & 0 & 1 & 0 & 0 & -1 & 0 & 0 & 0 & 0 & 0\\
0 & 0 & 0 & -1 & 0 & 0 & 0 & -1 & 0 & 0 & 1 & 0\\
0 & 1 & 0 & 0 & -1 & 0 & 0 & 0 & 0 & 0 & 0 & 0\\
0 & 1 & 1 & 0 & 0 & -1 & 0 & 0 & 1 & 0 & -1 & -1\\
\end{pmatrix},\qquad \mu=536.\]
\[d=(197,298,180,268,298,155,268,144,197,155,412,108).\]
\[B_0=\{0,1,2,4,5,7,8,9,10,11\};\quad B_1=\{3,6\}.\]
\endgroup\end{minipage}\par
\par\medskip\noindent\begin{minipage}{\linewidth}\textbf{F.5D.0569.}
\begingroup\small\setlength{\arraycolsep}{3pt}
\[U=\begin{pmatrix}
-1 & -1 & 0 & 1 & -1 & 0 & 0 & -1 & 1 & 0 & 0 & 1\\
0 & 0 & 0 & 0 & 0 & 0 & -1 & 0 & 0 & 0 & 0 & 1\\
0 & 1 & -1 & 0 & 0 & 0 & 0 & 0 & 0 & 0 & 0 & 0\\
0 & 0 & 0 & 0 & 1 & -1 & 0 & 0 & 0 & 0 & 0 & 0\\
1 & 1 & 0 & -1 & 0 & 0 & 0 & 0 & 0 & 1 & -1 & 0\\
\end{pmatrix},\qquad \mu=536.\]
\[d=(88,268,268,264,278,278,298,108,180,130,222,298).\]
\[B_0=\{0,3,4,5,6,7,8,9,10,11\};\quad B_1=\{1,2\}.\]
\endgroup\end{minipage}\par
\par\medskip\noindent\begin{minipage}{\linewidth}\textbf{F.5D.0841.}
\begingroup\small\setlength{\arraycolsep}{3pt}
\[U=\begin{pmatrix}
0 & -1 & 1 & 0 & 0 & 0 & 0\\
-1 & 0 & 1 & 0 & 0 & 0 & 0\\
0 & 0 & 1 & 0 & -1 & 0 & -1\\
0 & 0 & 0 & 0 & 0 & -1 & 1\\
0 & 0 & 1 & -1 & 0 & 0 & 0\\
\end{pmatrix},\qquad \mu=1250.\]
\[d=(1125,1125,1125,1125,500,625,625).\]
\[B_0=\{0,1,2,3,4\};\quad B_1=\{5,6\}.\]
\endgroup\end{minipage}\par
\par\medskip\noindent\begin{minipage}{\linewidth}\textbf{F.6D.1739.}
\begingroup\small\setlength{\arraycolsep}{3pt}
\[U=\begin{pmatrix}
0 & -1 & 0 & 0 & 0 & 0 & 0 & 0 & 1 & 0\\
0 & 0 & 0 & 0 & -1 & 0 & 0 & 1 & 0 & 0\\
-1 & 0 & 0 & 0 & 0 & 0 & 0 & 0 & 0 & 1\\
0 & 0 & 0 & 0 & 0 & 0 & -1 & 0 & 1 & 0\\
0 & 0 & 1 & 0 & 0 & -1 & 0 & 1 & 2 & -2\\
0 & 0 & 0 & -1 & 0 & 0 & 0 & 0 & 0 & 1\\
\end{pmatrix},\qquad \mu=7224.\]
\[d=(6048,4992,750,6048,3612,2250,4992,3612,4992,6048).\]
\[B_0=\{0,1,2,3,5,6,8,9\};\quad B_1=\{4,7\}.\]
\endgroup\end{minipage}\par
\par\medskip\noindent\begin{minipage}{\linewidth}\textbf{F.6D.2022.}
\begingroup\small\setlength{\arraycolsep}{3pt}
\[U=\begin{pmatrix}
0 & -1 & 0 & 0 & 0 & 0 & 0 & 0 & 1 & 0\\
0 & 0 & 0 & 1 & 0 & -1 & 0 & 0 & 0 & 0\\
-1 & 0 & 0 & 1 & 0 & 0 & 0 & 0 & 0 & 1\\
0 & 0 & 0 & 0 & 0 & 0 & 0 & -1 & 1 & 0\\
0 & 0 & 1 & -2 & 0 & 0 & -1 & 0 & 2 & -2\\
0 & 0 & 0 & 1 & -1 & 0 & 0 & 0 & 0 & 1\\
\end{pmatrix},\qquad \mu=7224.\]
\[d=(6372,5872,2000,3612,6372,3612,1000,5872,5872,2760).\]
\[B_0=\{0,1,2,4,6,7,8,9\};\quad B_1=\{3,5\}.\]
\endgroup\end{minipage}\par
\par\medskip\noindent\begin{minipage}{\linewidth}\textbf{F.6D.3068.}
\begingroup\small\setlength{\arraycolsep}{3pt}
\[U=\begin{pmatrix}
0 & 0 & 0 & -1 & 0 & 0 & 0 & 0 & 0 & 0 & 1 & 0\\
0 & 0 & 0 & 0 & 0 & 0 & 0 & 1 & 0 & -1 & 0 & 0\\
0 & -1 & 0 & 0 & 0 & 0 & 0 & 0 & 0 & 0 & 0 & 1\\
1 & 0 & 0 & 0 & -1 & 0 & 0 & 0 & 0 & 0 & -1 & 1\\
-1 & 0 & -1 & 0 & 0 & 0 & 1 & 0 & 0 & 0 & 1 & 0\\
0 & 0 & 0 & 0 & 0 & -1 & 0 & 1 & 1 & 0 & 1 & -1\\
\end{pmatrix},\qquad \mu=5832.\]
\[d=(2447,3813,2800,3287,2973,3585,1960,2916,1195,2916,3287,3813).\]
\[B_0=\{0,1,2,3,4,5,6,8,10,11\};\quad B_1=\{7,9\}.\]
\endgroup\end{minipage}\par
\par\medskip\noindent\begin{minipage}{\linewidth}\textbf{F.6D.3115.}
\begingroup\small\setlength{\arraycolsep}{3pt}
\[U=\begin{pmatrix}
-1 & 0 & 0 & 0 & 0 & 1 & 0 & 1 & 0 & 0 & -1 & 0\\
-1 & 0 & 0 & -1 & 0 & 0 & 0 & 1 & 0 & 0 & 0 & 1\\
0 & 0 & 0 & 0 & 1 & 0 & 0 & 0 & 0 & -1 & 0 & 0\\
0 & 1 & 0 & 0 & 0 & 0 & 0 & 0 & -1 & 0 & 0 & 0\\
1 & 1 & -1 & 0 & 0 & 0 & 0 & -1 & 0 & 0 & 0 & 0\\
1 & 1 & 0 & 0 & -1 & 0 & -1 & -1 & 0 & 0 & 0 & 0\\
\end{pmatrix},\qquad \mu=5848.\]
\[d=(2773,4058,4764,2067,2924,2773,1840,2067,4058,2924,2067,2773).\]
\[B_0=\{0,1,2,3,5,6,7,8,10,11\};\quad B_1=\{4,9\}.\]
\endgroup\end{minipage}\par
\par\medskip\noindent\begin{minipage}{\linewidth}\textbf{F.6D.3127.}
\begingroup\small\setlength{\arraycolsep}{3pt}
\[U=\begin{pmatrix}
0 & 0 & 0 & 0 & 0 & 1 & -1 & 0 & 0 & 0 & 0 & 0\\
0 & -1 & 1 & 0 & 0 & 0 & 0 & 0 & 0 & 0 & 0 & 0\\
0 & 0 & 0 & 0 & -1 & 0 & 0 & 0 & 0 & 0 & 0 & 1\\
0 & 0 & 0 & 0 & 0 & -1 & 0 & 0 & 0 & -1 & 1 & 1\\
0 & 0 & 0 & 0 & 0 & 1 & 0 & 1 & -1 & 0 & 0 & -1\\
1 & 0 & 1 & -1 & 0 & 1 & 0 & 0 & 0 & 0 & 0 & -1\\
\end{pmatrix},\qquad \mu=5848.\]
\[d=(1210,2924,2924,3630,3932,3428,3428,2672,2168,2672,2168,3932).\]
\[B_0=\{0,3,4,5,6,7,8,9,10,11\};\quad B_1=\{1,2\}.\]
\endgroup\end{minipage}\par
\par\medskip\noindent\begin{minipage}{\linewidth}\textbf{F.6D.3996.}
\begingroup\small\setlength{\arraycolsep}{3pt}
\[U=\begin{pmatrix}
0 & 0 & 0 & -1 & 0 & -1 & 1 & 0 & 0 & 0 & 0\\
0 & 0 & 0 & 0 & 0 & 0 & 0 & 0 & -1 & 1 & 0\\
0 & 0 & -1 & 0 & 0 & 0 & 1 & 0 & 0 & 0 & 0\\
0 & 1 & 0 & 0 & -1 & 0 & 0 & 0 & 0 & 0 & 0\\
-1 & 1 & 0 & 0 & 0 & 1 & -1 & 0 & 0 & 0 & 0\\
0 & 1 & 0 & 0 & 0 & 0 & -1 & -1 & 0 & 1 & 1\\
\end{pmatrix},\qquad \mu=7404.\]
\[d=(2504,5304,5896,2800,5304,3096,5896,4665,3702,3702,1555).\]
\[B_0=\{0,1,2,3,4,5,6,7,10\};\quad B_1=\{8,9\}.\]
\endgroup\end{minipage}\par
\par\medskip\noindent\begin{minipage}{\linewidth}\textbf{F.6D.4188.}
\begingroup\small\setlength{\arraycolsep}{3pt}
\[U=\begin{pmatrix}
0 & -1 & 0 & 0 & 1 & 0 & -1 & 0 & 0 & 0 & 0 & 0\\
0 & 0 & 0 & 0 & 0 & 1 & 0 & 0 & 0 & 0 & 0 & -1\\
-1 & 0 & 0 & 0 & 0 & -1 & 1 & 0 & -1 & 1 & 0 & 0\\
0 & 0 & 1 & 0 & 0 & 0 & 0 & 0 & 0 & 0 & -1 & 0\\
1 & 0 & 0 & 0 & 1 & 0 & -1 & -1 & 0 & 0 & 0 & 0\\
0 & 0 & 1 & -1 & -1 & 0 & 1 & 0 & 0 & 0 & 0 & 0\\
\end{pmatrix},\qquad \mu=6376.\]
\[d=(1871,2020,4660,2640,4529,3188,2509,3891,1275,3825,4660,3188).\]
\[B_0=\{0,1,2,3,4,6,7,8,9,10\};\quad B_1=\{5,11\}.\]
\endgroup\end{minipage}\par
\par\medskip\noindent\begin{minipage}{\linewidth}\textbf{F.6D.4954.}
\begingroup\small\setlength{\arraycolsep}{3pt}
\[U=\begin{pmatrix}
0 & 1 & 0 & -1 & -1 & -1 & 0 & 0 & 0 & 0 & 1\\
0 & 0 & 0 & 1 & 0 & 0 & -1 & 0 & 0 & 1 & 0\\
0 & 0 & 0 & 0 & 0 & 0 & 0 & -1 & 0 & 0 & 1\\
0 & 0 & -1 & 0 & 1 & 0 & 0 & 0 & 0 & 0 & 0\\
0 & 0 & 0 & 1 & 0 & 0 & 0 & 0 & -1 & 0 & 0\\
-1 & 0 & 0 & 0 & 0 & 0 & 0 & 0 & 0 & 1 & 1\\
\end{pmatrix},\qquad \mu=7452.\]
\[d=(6453,4680,3726,3887,3726,1560,5847,4493,3887,1960,4493).\]
\[B_0=\{0,1,3,5,6,7,8,9,10\};\quad B_1=\{2,4\}.\]
\endgroup\end{minipage}\par
\par\medskip\noindent\begin{minipage}{\linewidth}\textbf{F.6D.5485.}
\begingroup\small\setlength{\arraycolsep}{3pt}
\[U=\begin{pmatrix}
0 & 0 & 0 & 0 & 0 & 0 & 0 & -1 & 0 & 0 & 0 & 0 & 1 & 0\\
1 & 0 & 1 & 0 & 0 & 0 & 0 & 0 & -1 & -1 & 0 & -1 & -1 & 1\\
0 & 0 & 0 & 0 & 0 & 0 & 0 & 0 & 0 & 0 & -1 & 0 & 0 & 1\\
0 & 0 & 0 & 0 & -1 & 0 & 0 & 0 & 1 & 0 & 0 & 0 & 0 & 0\\
0 & -1 & -1 & 1 & 0 & -1 & 0 & 0 & 1 & 0 & 0 & 1 & 1 & 0\\
0 & 1 & 0 & 0 & 0 & 0 & -1 & 0 & 0 & 0 & 0 & 0 & 0 & 0\\
\end{pmatrix},\qquad \mu=5160.\]
\[d=(1850,2880,2340,1170,2580,1850,2880,2580,2580,1170,2880,740,2580,2880).\]
\[B_0=\{0,1,2,3,5,6,9,10,11,13\};\quad B_1=\{4,8\};\quad B_2=\{7,12\}.\]
\endgroup\end{minipage}\par
\par\medskip\noindent\begin{minipage}{\linewidth}\textbf{F.6D.5541.}
\begingroup\small\setlength{\arraycolsep}{3pt}
\[U=\begin{pmatrix}
0 & 0 & 0 & 0 & 0 & 0 & -1 & 0 & 1 & 0 & 0 & 0\\
0 & 0 & -1 & 0 & 1 & 0 & 0 & 0 & 0 & 0 & 0 & 0\\
0 & 0 & 0 & -1 & 0 & 0 & 0 & 0 & 0 & 0 & 0 & 1\\
0 & 1 & 0 & 0 & 0 & 0 & 0 & 0 & 0 & 0 & -1 & 0\\
1 & 0 & 0 & 0 & 0 & 0 & 0 & 0 & 0 & -1 & 0 & 0\\
1 & -1 & 0 & 0 & 1 & -1 & 0 & 1 & 1 & 0 & 0 & -1\\
\end{pmatrix},\qquad \mu=6496.\]
\[d=(3248,3792,3248,3792,3248,3240,3248,1080,3248,3248,3792,3792).\]
\[B_0=\{1,3,5,7,10,11\};\quad B_1=\{0,9\};\quad B_2=\{2,4\};\quad B_3=\{6,8\}.\]
\endgroup\end{minipage}\par
\par\medskip\noindent\begin{minipage}{\linewidth}\textbf{F.6D.5880.}
\begingroup\small\setlength{\arraycolsep}{3pt}
\[U=\begin{pmatrix}
1 & 0 & 0 & 0 & -1 & -1 & 0 & 0 & 0 & -1 & 0 & 0 & 1 & 1\\
0 & 1 & 0 & 0 & 0 & 0 & 0 & 0 & -1 & 0 & 0 & 0 & 0 & 0\\
0 & -1 & 0 & 0 & 0 & 0 & 0 & -1 & 0 & 0 & 0 & 0 & 0 & 1\\
0 & 0 & 0 & 0 & 0 & 1 & -1 & 0 & 0 & 0 & 0 & 0 & 0 & 0\\
0 & 0 & 0 & -1 & 0 & 1 & 0 & 0 & 0 & 1 & -1 & 1 & -1 & 0\\
0 & 0 & -1 & 1 & 0 & 0 & 0 & 0 & 0 & 0 & 0 & 0 & 0 & 0\\
\end{pmatrix},\qquad \mu=5160.\]
\[d=(1110,2580,2830,2830,1950,2830,2830,1390,2580,1665,1665,1365,1365,3970).\]
\[B_0=\{0,2,3,4,5,6,7,9,10,11,12,13\};\quad B_1=\{1,8\}.\]
\endgroup\end{minipage}\par
\par\medskip\noindent\begin{minipage}{\linewidth}\textbf{F.6D.5962.}
\begingroup\small\setlength{\arraycolsep}{3pt}
\[U=\begin{pmatrix}
0 & 0 & 0 & 0 & 0 & -1 & 0 & 0 & 0 & 1 & 0 & 0\\
0 & 0 & -1 & 0 & 0 & 0 & 0 & 0 & 0 & 0 & 0 & 1\\
1 & -1 & 0 & 0 & 0 & 0 & 0 & 0 & 0 & 0 & 0 & 0\\
0 & 0 & 0 & 0 & 0 & 0 & -1 & 0 & 1 & 0 & 0 & -1\\
0 & 0 & 0 & 0 & -1 & 0 & 0 & 0 & 0 & 0 & 1 & 0\\
-1 & 0 & 0 & 1 & 0 & 0 & 0 & -1 & -1 & 1 & 1 & 0\\
\end{pmatrix},\qquad \mu=6496.\]
\[d=(3384,3384,3248,2700,3656,3656,1760,1620,5008,3656,3656,3248).\]
\[B_0=\{0,1,3,4,5,6,7,8,9,10\};\quad B_1=\{2,11\}.\]
\endgroup\end{minipage}\par
\par\medskip\noindent\begin{minipage}{\linewidth}\textbf{F.6D.5965.}
\begingroup\small\setlength{\arraycolsep}{3pt}
\[U=\begin{pmatrix}
0 & 0 & 0 & 0 & -1 & 0 & 0 & 0 & 0 & 0 & 1\\
0 & -1 & 0 & 0 & 0 & 0 & 0 & 0 & 0 & 1 & 0\\
0 & 0 & 0 & 0 & 0 & -1 & 0 & 1 & 0 & 0 & 0\\
0 & 0 & -1 & 0 & 0 & 0 & -1 & 0 & 1 & 1 & -1\\
0 & 0 & 1 & -1 & 0 & 0 & 0 & 0 & 0 & 0 & 0\\
-1 & 0 & 0 & 0 & 0 & 0 & 0 & -1 & 1 & 1 & 0\\
\end{pmatrix},\qquad \mu=7236.\]
\[d=(4950,4266,3618,3618,3618,4266,1980,4266,4950,4266,3618).\]
\[B_0=\{0,1,5,6,7,8,9\};\quad B_1=\{2,3\};\quad B_2=\{4,10\}.\]
\endgroup\end{minipage}\par
\par\medskip\noindent\begin{minipage}{\linewidth}\textbf{F.6D.6078.}
\begingroup\small\setlength{\arraycolsep}{3pt}
\[U=\begin{pmatrix}
0 & -1 & 0 & 0 & 0 & -1 & -1 & 0 & 0 & -1 & 1 & 0 & 1 & 1\\
0 & 0 & 0 & 0 & 0 & 0 & 0 & 0 & -1 & 0 & 0 & 1 & 0 & 0\\
-1 & 0 & 0 & 0 & 0 & 0 & 0 & 0 & 0 & 0 & 0 & 0 & 0 & 1\\
0 & 0 & 0 & 0 & -1 & 1 & 0 & 0 & 0 & 0 & 0 & 0 & 0 & 0\\
0 & 1 & -1 & 0 & 0 & 0 & 0 & 0 & 0 & 0 & 0 & 0 & 0 & 0\\
0 & 0 & 0 & 1 & 0 & 1 & 0 & -1 & 0 & 1 & -1 & 1 & 0 & -1\\
\end{pmatrix},\qquad \mu=5160.\]
\[d=(2980,2680,2680,1110,2580,2580,1110,1950,2680,780,2220,2680,1950,2980).\]
\[B_0=\{0,1,2,3,6,7,8,9,10,11,12,13\};\quad B_1=\{4,5\}.\]
\endgroup\end{minipage}\par
\par\medskip\noindent\begin{minipage}{\linewidth}\textbf{F.6D.6547.}
\begingroup\small\setlength{\arraycolsep}{3pt}
\[U=\begin{pmatrix}
0 & 0 & 0 & 1 & 0 & 0 & -1 & 0 & 0 & 0\\
0 & -1 & 0 & 0 & 0 & 0 & 0 & 0 & 0 & 1\\
0 & 0 & -1 & 0 & 0 & 0 & 0 & 0 & 0 & 1\\
1 & 0 & 0 & 0 & -1 & 0 & 0 & 0 & 0 & 0\\
1 & 0 & 0 & 1 & 0 & 0 & 0 & -1 & 1 & -2\\
0 & 0 & 0 & 0 & 0 & -1 & 0 & 0 & 0 & 1\\
\end{pmatrix},\qquad \mu=\frac{25984}{3}.\]
\[d=(5056,6496,6496,5056,5056,6496,5056,1440,4320,6496).\]
\[B_0=\{0,3,4,6,7,8\};\quad B_1=\{1,2,5,9\}.\]
\endgroup\end{minipage}\par
\par\medskip\noindent\begin{minipage}{\linewidth}\textbf{F.6D.6614.}
\begingroup\small\setlength{\arraycolsep}{3pt}
\[U=\begin{pmatrix}
0 & 0 & 0 & 0 & 0 & 0 & 1 & -1 & 0 & 0 & 0\\
0 & 0 & 0 & 0 & 0 & -1 & 0 & 0 & 0 & 0 & 1\\
-1 & 0 & 1 & 1 & 0 & 0 & 0 & 0 & 0 & 1 & -1\\
0 & -1 & 1 & 1 & 0 & 0 & -1 & 0 & 0 & 1 & 0\\
0 & 0 & 0 & 0 & -1 & 0 & 0 & 0 & 0 & 1 & 0\\
0 & 0 & 1 & 0 & 0 & 0 & 0 & 0 & -1 & 1 & 0\\
\end{pmatrix},\qquad \mu=7236.\]
\[d=(4455,4455,2025,2970,3618,4158,4158,4158,5643,3618,4158).\]
\[B_0=\{0,1,2,3,5,6,7,8,10\};\quad B_1=\{4,9\}.\]
\endgroup\end{minipage}\par
\par\medskip\noindent\begin{minipage}{\linewidth}\textbf{F.6D.7472.}
\begingroup\small\setlength{\arraycolsep}{3pt}
\[U=\begin{pmatrix}
0 & -1 & 0 & 0 & 0 & 0 & 0 & 0 & 1 & 0\\
-1 & 0 & 0 & 1 & 0 & 0 & 0 & 0 & 0 & 1\\
0 & 0 & 0 & 1 & 0 & -1 & 0 & 0 & 0 & 0\\
0 & 0 & 0 & 0 & 0 & 0 & 0 & -1 & 1 & 0\\
0 & 0 & 1 & -1 & 0 & 0 & -1 & 0 & 1 & -1\\
0 & 0 & 0 & 1 & -1 & 0 & 0 & 0 & 0 & 1\\
\end{pmatrix},\qquad \mu=9024.\]
\[d=(7632,6352,4480,4512,7632,4512,3200,6352,6352,3120).\]
\[B_0=\{0,1,2,4,6,7,8,9\};\quad B_1=\{3,5\}.\]
\endgroup\end{minipage}\par
\par\medskip\noindent\begin{minipage}{\linewidth}\textbf{F.6D.7551.}
\begingroup\small\setlength{\arraycolsep}{3pt}
\[U=\begin{pmatrix}
0 & -1 & 1 & 0 & 0 & 0 & 0 & 0\\
-1 & 0 & 1 & 0 & 0 & 0 & 0 & 0\\
0 & 0 & 1 & -1 & 0 & 0 & 0 & 0\\
0 & 0 & 1 & 0 & 0 & -1 & 0 & -1\\
0 & 0 & 0 & 0 & 0 & 0 & -1 & 1\\
0 & 0 & 1 & 0 & -1 & 0 & 0 & 0\\
\end{pmatrix},\qquad \mu=15552.\]
\[d=(14256,14256,14256,14256,14256,6480,7776,7776).\]
\[B_0=\{0,1,2,3,4,5\};\quad B_1=\{6,7\}.\]
\endgroup\end{minipage}\par


\begin{thebibliography}{BDGP21}
\bibitem[BDGP21]{BDGP} I. Biswas, A. Dey, O. Gen\c c, and M. Poddar,
\emph{On stability of tangent bundle of toric varieties},
Proc. Indian Acad. Sci. Math. Sci. \textbf{131} (2021), Paper No.~36,
arXiv:1808.08701.
\bibitem[CLS11]{CLS} D. Cox, J. Little, and H. Schenck,
\emph{Toric Varieties}, Graduate Studies in Mathematics, vol.~124,
American Mathematical Society, 2011.
\bibitem[CT24]{CT} A. Clarke and C. Tipler,
\emph{Semi-stability and local wall-crossing for Hermitian Yang--Mills
connections}, J. Symplectic Geom. \textbf{22} (2024), no.~3,
673--694; arXiv:2304.05245.
\bibitem[DDK20]{DDK} J. Dasgupta, A. Dey, and B. Khan,
\emph{Stability of equivariant vector bundles over toric varieties},
Doc. Math. \textbf{25} (2020), 1787--1833.
\bibitem[Fuj84]{Fujishige} S. Fujishige,
\emph{A characterization of faces of the base polyhedron associated
with a submodular system}, J. Oper. Res. Soc. Japan \textbf{27} (1984),
no.~2, 112--129,
\href{https://doi.org/10.15807/jorsj.27.112}{doi:10.15807/jorsj.27.112}.
\bibitem[GRT16]{GRT} D. Greb, J. Ross, and M. Toma,
\emph{Variation of Gieseker moduli spaces via quiver GIT},
Geom. Topol. \textbf{20} (2016), 1539--1610.
\bibitem[HNS22]{HNS} M. Hering, B. Nill, and H. S\"u\ss,
\emph{Stability of tangent bundles on smooth toric Picard-rank-2
varieties and surfaces}, Facets of Algebraic Geometry, Vol.~II,
London Math. Soc. Lecture Note Ser., vol.~473, Cambridge University
Press, 2022, 1--25; arXiv:1910.08848.
\bibitem[Kly90]{Klyachko} A. A. Klyachko,
\emph{Equivariant bundles on toral varieties}, Math. USSR-Izv.
\textbf{35} (1990), 337--375.
\bibitem[NT24]{NT} A. Napame and C. Tipler,
\emph{Toric sheaves, stability and fibrations}, J. Pure Appl. Algebra
\textbf{228} (2024), Paper No.~107494; arXiv:2210.04587.
\bibitem[Rey23]{Reynolds} W. Reynolds,
\emph{On slope stability of tangent sheaves on smooth toric Fano
varieties}, Ph.D. thesis, University of Edinburgh, 2023,
\href{https://doi.org/10.7488/era/3511}{doi:10.7488/era/3511}.
\bibitem[Wue26]{Wuebben} B. J. Wuebben,
\emph{Slope stability of tangent bundles of smooth toric Fano varieties},
arXiv:2608.20411, 2026.
\end{thebibliography}
\end{document}
